\section{Discussion}

\subsection{What CIKD++-RT Learned}
Our analysis of the CIKD++-RT framework demonstrates its ability to effectively synthesize multimodal signals (Text, Image, and Knowledge Graph) for fine-grained multimodal misinformation detection. The model successfully learned how to leverage the TVCS module to explicitly model interactions between textual claims and external knowledge entities \cite{ref43}. This is directly reflected in the raw TVCS score distribution metrics: the mean score of the Real class is kept low (0.304671), while the CK class is pushed to a high level (0.503780), creating a significant Delta separation (0.199109) and achieving an impressive AUC of 0.726705. The core essence that CIKD++-RT internalized during training is knowledge-grounded structural contradiction signals, rather than absolute truth understanding or real-world fact-checking capabilities \cite{ref42}.

\subsection{Why no\_c\_emb Outperforms Alternative Methods}
During architecture development and refinement, the no\_c\_emb variant (which removes the scalar class embedding component c\_emb) proved to be the optimal configuration. Internal diagnostic tests (specifically the G3 test) indicate that the learned weight (beta) of c\_emb remained almost dormant at its initial initialization level throughout the training process. This mathematically demonstrates that c\_emb provided no significant feature extraction value. Instead, it acted as a noisy "shortcut," causing the neural network to easily fall into overfitting based on class priors rather than actual data cross-referencing \cite{ref44}.

By entirely eliminating this component, the architecture forced the model to abandon its "guessing" tendency. The neural network is now compelled to rely entirely on the cross-attention mechanism and the high-quality multimodal visual evidence vector ($z_v$) synthesized from the TVCS module. This architectural streamlining ensures that the cross-referencing between Text, Image, and knowledge graph occurs more rigorously, preventing information noise, thereby directly improving the model's generalization capability on the locked test split (enabling the best variant to achieve a Macro-F1 of 0.4792 and a CK-F1 of 0.3922 during the diagnostic phase).

\subsection{Improvements in CK Label: Increased Precision and Reduced False Positives}
A critical observation regarding the performance of the content-knowledge inconsistency (CK) class is that the improvements were entirely driven by an increase in precision (precision increased 31.93\% $\rightarrow$ 38.74\%) and a reduction in false-positive predictions (false-positive decreased 194 $\rightarrow$ 136), rather than an increase in recall. In fact, we observed that the CK recall slightly dropped from 38.56\% to 36.44\%. "F4 is not a comprehensively stronger classifier by a large margin. Its most apparent contribution is the precision improvement focused on the CK label and the Real-vs-CK evidence separation capability of the TVCS module. The increase in the Macro-F1 metric is positive but modest; the recall of the CK label and the performance on class 5 remain persistent limitations."

\subsection{Calibration}
The initial raw F4 model exhibited extreme overconfidence (raw ECE of 14.80\%, higher than the baseline). However, by applying the post-hoc temperature scaling calibration method (T=1.522523) \cite{ref10}, we successfully reduced the ECE on the locked test split to 3.85\% without altering any predictions (zero argmax changes). The model produces confidence scores that more closely track empirical probabilities for specific predictions. We note that this does not imply the raw calibration strictly outperforms all baselines across every metric, but rather that the architectural refinements help mitigate extreme overconfidence \cite{ref45} in misclassified samples.

\subsection{Limitations on SOTA (State-of-the-Art) Comparisons}
Contextualizing our findings against current State-of-the-Art (SOTA) methods is crucial. We refrain from making any claims that CIKD++-RT achieved absolute direct SOTA performance because we have not reproduced systems like KGAlign \cite{ref46}/KAMP \cite{ref47} on the same-split dataset; instead, we demonstrate the internal architectural superiority of our method against the passive Text+Image+KG baseline, specifically for the fine-grained contradiction detection task.

\subsection{Future Work}
Future iterations of this research need to address the remaining challenges in multimodal misinformation. The primary focus should be algorithmic strategies to boost CK recall without sacrificing the newly attained precision \cite{ref48}. Furthermore, resolving the difficulty in classifying minority classes (e.g., class 5) \cite{ref49} and exploring more sophisticated knowledge graph integrations will be vital steps toward a more comprehensive contradiction detection framework. These experiments will be conducted on external holdout sets, with absolutely no further tuning on the current locked test split.
